% =============================================================================
\section{Cross-Dimensional Interactions: Lessons from Our Work}
\label{sec:cross}
% =============================================================================

Cross-dimensional interactions are easy to miss when each dimension
is studied separately, but they dominate real-world performance.
Below we formalize the main interactions our publications have
identified.


% -----------------------------------------------------------------------------
\subsection{Workload $\times$ Router}
\label{sec:cross:wr}
% -----------------------------------------------------------------------------

\begin{observation}[Workload determines router adaptation mechanism]
\label{obs:wr-adapt}
Chat workloads produce per-turn user feedback; our Feedback
Detector~\cite{feedbackdet2026} classifies satisfaction in real
time, allowing the router to re-route on dissatisfaction without
explicit reward modeling.  Agent workloads lack such per-turn
signals---task success is observed only at the end of a multi-step
session---so RL-based routing~\cite{routerr1_2025} that reasons
over sequential decisions is a better fit.  The adaptation
mechanism must match the workload's feedback granularity.
\end{observation}

\begin{observation}[Modality determines routing signals]
\label{obs:wr-modal}
Text-only chat uses embedding similarity and keyword
signals~\cite{vllmsr2026}.  VLM workloads require multimodal
classifiers---SigLIP+MiniLM in AVR~\cite{avr2026}---because
text-only signals cannot assess screenshot complexity.  The router signal set must match the workload modality.
\end{observation}

\begin{observation}[Workload determines routing safety and privacy requirements]
\label{obs:wr-safety}
Different workload types impose different safety and privacy
requirements on the router.  Chat routing errors produce quality
degradation; multi-turn chat adds a subtler threat---adversarial
users can spread jailbreak content across turns, evading
single-turn classifiers.  The \vsr{}~\cite{vllmsr2026} addresses
this with contrastive embedding classification that evaluates turn
sequences, not individual messages.  CUA routing errors produce
outright \emph{security vulnerabilities}: the Visual Confused
Deputy~\cite{vcd2026} shows that misrouting a CUA action to a
less-capable model can cause grounding errors exploitable for
privilege escalation.  Across all workload types, our hierarchical
MLCommons safety pipeline---binary gate~\cite{safetyl1_2026} then
nine-class hazard categorization~\cite{safetyl2_2026}---provides
uniform content filtering at the routing layer, with category-level
granularity that lets operators enforce workload-specific policies
(e.g., blocking privacy-exfiltrating prompts in customer-facing chat,
flagging misinformation in knowledge-retrieval agents).  Safety and privacy should be integrated as routing objectives,
with multi-turn awareness for chat and action-level guardrails for
agents.
\end{observation}

\begin{observation}[Factual workloads require response-time validation]
\label{obs:wr-halu}
Request-time routing alone cannot guarantee correctness: a model may
be well-suited for the query's domain yet still hallucinate on a
specific fact.  The Factcheck Classifier~\cite{factcheck2026}
identifies fact-seeking prompts at request time ($\sim$12\,ms), and
HaluGate~\cite{halugate2025} validates responses at token level
(76--162\,ms overhead).  Together they close a loop that
request-time routing cannot: per-model hallucination rates observed
by HaluGate feed back into the routing policy, gradually steering
factual traffic toward more reliable models.
\end{observation}


% -----------------------------------------------------------------------------
\subsection{Router $\times$ Pool}
\label{sec:cross:rp}
% -----------------------------------------------------------------------------

\begin{observation}[Router policy co-determines pool sizing]
\label{obs:rp-codesign}
FleetOpt~\cite{fleetopt2026} shows that co-designing the compression
parameter $\gamma$ with pool sizes $(n_s, n_l)$ yields 3.1--6.4\%
lower cost than retrofitting compression onto a pre-existing fleet.
The router is a co-design variable, not a layer atop the pool.
\end{observation}

\begin{observation}[Router latency constrains pool topology]
\label{obs:rp-latency}
FastRouter~\cite{fastrouter2026} showed that routing latency is
a critical constraint: at 4,918\,ms baseline, routing overhead
exceeds the TTFT SLO for short-context requests.  Only after the
98$\times$ reduction to 50\,ms does pool routing become practical
for the full request stream.
mmBERT-Embed~\cite{mmberted2025} reinforces this: its 2D Matryoshka
design lets operators trade embedding quality for latency
(2.56\,ms at layer-6 vs.\ 11\,ms at layer-22), keeping the
similarity signal within the router's latency budget even as pool
topology grows more complex.
\end{observation}

\begin{observation}[Mid-session model switching incurs a KV-cache sunk cost]
\label{obs:rp-sunkcost}
In multi-turn sessions, if the Feedback
Detector~\cite{feedbackdet2026} or context memory signals indicate
that a better model is available, the router faces a dilemma.
Switching models mid-session discards the KV-cache accumulated for
the current model---a \emph{sunk cost} proportional to session
length and context size.  Not switching preserves the cache
investment but forgoes the quality improvement---an \emph{opportunity
cost} that grows with every subsequent turn served by the
suboptimal model.  The optimal policy depends on the remaining
session length, the quality gap between models, and the KV-cache
rebuild cost, creating a Pareto frontier between sunk cost and opportunity cost.
No current system explicitly optimizes this trade-off.
\end{observation}


% -----------------------------------------------------------------------------
\subsection{Workload $\times$ Pool}
\label{sec:cross:wp}
% -----------------------------------------------------------------------------

\begin{observation}[Workload archetype determines optimal topology]
\label{obs:wp-topo}
Our fleet tools~\cite{fleetopt2026, fleetsim2026} show:
Archetype~1 (concentrated-below) benefits most from two-pool routing.
Archetype~3 (concentrated-above) may benefit from disaggregated
prefill/decode because most requests are long.  Archetype~2
(dispersed) may require multi-tier pooling.
\end{observation}

\begin{observation}[Agent adoption shifts the energy frontier]
\label{obs:wp-energy}
The 1/W law~\cite{onewlaw2026} predicts that shifting from chat
to agents drops fleet-wide energy efficiency by $2$--$4\times$
unless pool topology adapts.  Two-pool routing recovers most of
this loss, but only if pool boundaries track workload evolution.
\end{observation}


